\subsection{Future Work}
\subsubsection{Deployment on Real Underwater Vehicles and Full SLAM Integration}
The most important future work is to integrate and test the complete Side Scan Sonar SLAM data processing pipeline on real underwater vehicles with side scan sonar capability. The results in this thesis show that the pipeline can process real sensor data in real time on embedded hardware, but the next step is to run it as part of a complete onboard autonomy stack during real missions. Suitable platforms include AUR Lab LAUV vehicles equipped with side scan sonar, BlueROV based platforms with external sonar payloads, or other autonomous underwater vehicles that already support navigation sensors and side scan sonar logging.
\\ \\
The first goal should be practical deployment rather than further theoretical extension. This means installing the ROS2 pipeline on the vehicle, connecting it to the real sonar, IMU, DVL, depth sensor, and navigation topics, and validating that the generated local maps and landmarks remain stable during actual underwater operation. Since the software stack has been made modular, configurable, and portable, this integration should mainly require platform specific configuration of topics, sensor frames, sonar parameters, and timing. For ROS2 based platforms, the existing packages can be used directly. For vehicles that do not use ROS2, the core libraries can still be reused without the ROS2 wrappers.
\\ \\
After the pipeline has been deployed and validated on a real vehicle, the next major step is completing the full closed loop SLAM system. The current work provides the real-time local map generation, feature extraction, landmark descriptors, and software infrastructure needed for this. Future work should therefore focus on data association, loop closure, and factor graph optimization. Landmark descriptors can be used for fast candidate pruning, while geometric consistency methods such as JCBB can be used to reject false associations. These landmark observations and loop closures can then be inserted as factors into an incremental optimizer such as iSAM2 in GTSAM.
\\ \\
Completing this loop would allow the system to correct dead reckoning drift and generate more globally consistent maps. The local maps produced by this pipeline are already useful for extracting stable features, but without loop closure the global map will still inherit long term state estimation drift. A complete SLAM backend would allow landmark matches to correct the vehicle trajectory, which would then improve future map generation and feature extraction. This creates the intended feedback loop between mapping, feature extraction, data association, and state estimation. Therefore, the most important future research direction is to take the implemented real-time data processing pipeline, deploy it on a real drone, and extend it into a complete closed loop Side Scan Sonar SLAM system.

\newpage

\subsubsection{Pipeline Improvements, Generalization, and Further Optimization}
A second important direction is improving the flexibility and performance of the processing library further. The current software stack is already modular, configurable, and separated into reusable processing libraries with ROS2 wrappers on top. This makes the pipeline portable and easier to integrate on different vehicles. However, the current implementation is still mainly built around side scan sonar geometry. Future work could extend this modular structure with additional sonar models, such as front-facing sonar, imaging sonar, multibeam sonar, or other acoustic sensors. The main change would be the sonar model and generated measurement geometry, while much of the existing map generation, feature extraction, and landmark descriptor framework could be reused.
\\ \\
This is one of the strengths of the current software design. Since the pipeline was implemented with scalability, adaptability, documentation, and portability in mind, adding support for new sonar profiles should not require rewriting the full system. Each sonar type could define its own beam geometry, projection model, measurement uncertainty, and valid return region, while the rest of the pipeline continues to use the same mapping and landmark extraction structure. This would make the already existing library foundation more general and useful across a wider range of underwater robots and acoustic perception systems.
\\ \\
There is also clear room for further performance optimization. The results show that memory usage and map resolution have a strong effect on runtime. Future work should therefore investigate better memory handling, compressed map representations, and more efficient data transport between processing stages. Lossless compression may be useful for reducing the amount of map data moved through memory and published between nodes, while still preserving the original information. This could be especially valuable for dense local maps, where a large part of the runtime cost comes from moving and publishing image like data.
\\ \\
The local map generation algorithm could also be accelerated using parallel processing. Many parts of the algorithm are similar to rasterization and illumination problems in classical computer graphics, where many pixels or samples can be processed independently. This makes the pipeline a good candidate for GPU acceleration. A GPU implementation could allow larger maps, higher resolution sonar data, or more detailed probabilistic models to run within the same embedded compute budget. Even small single board computers often include integrated GPUs, so this could make high fidelity sonar mapping possible without consuming several CPU cores.
\\ \\
Feature extraction can also be improved. The current implementation works in real time, but profiling shows that height estimation, filtering, and segmentation dominate much of the runtime. Future work should improve the landmark height estimation model, since the current method is only a rough approximation based on sonar shadow geometry and assumes a mostly flat seabed, while ridges, slopes, rocks, and uneven terrain can distort shadows and affect the estimated object height. Improved shadow modelling, better shadow segmentation, and more efficient algorithm and data structure design could make height estimation more reliable, less sensitive to noise, and faster to compute. One possible improvement is to fuse the detected shadow of each landmark with the pitch and roll angles of the sonar at the time the landmark was observed. This is challenging, since each landmark may cover several map cells and sonar swaths, making it non trivial to decide which attitude measurement should be used. However, smarter data structures that preserve the link between landmarks, contributing swaths, timestamps, and vehicle attitude could make this possible without adding too much runtime overhead. Terrain information, such as a local elevation estimate or seabed height map, could also compensate for shadow distortion on non flat terrain and improve height estimates. These methods must still run in real time, where GPU acceleration, shader style processing, and computer graphics inspired techniques could help keep the added modelling cost low.
\\ \\
Finally, the landmark representation itself can be improved. Nearby or overlapping landmark bounding boxes could be merged to reduce the number of features and compress the global landmark database further. Better descriptor pruning could make loop closure search faster and more reliable. Together, these improvements would make the pipeline more efficient, more general, and easier to deploy across different underwater vehicles and sonar configurations. The current work provides a strong foundation, but future work can expand it from a real-time side scan sonar pipeline into a broader acoustic perception framework for underwater autonomy.
